% !TEX root = tarasov-qasp-zh.tex
\chapter[考虑摩擦时动力学问题会难解得多吗？]{考虑摩擦时，动力学问题\\会难解得多吗？\label{ch-07}}

\begin{dialogue}
\tchr 把摩擦力考虑进来，习题可能变得难解得多。

\std 可是摩擦力我们已经讨论过了（\chap{ch-03}）。物体运动时，摩擦力由支承力确定（$F_{f\!\!r}=kN$）；物体静止时，摩擦力等于力图使它脱离这一静止状态的那个力。这些都容易理解，也容易记住。

\tchr 话是不错。不过你忽略了一个重要事实：你已经默认了下列问题的答案是已知的——
\begin{enumerate*}[label=(\arabic*)]
\item 物体是在运动，还是在静止？
\item 物体（如果确实在运动）朝哪个方向运动？
\end{enumerate*}

这些若事先已知，问题就比较简单。否则，问题可能从一开始就非常复杂，甚至需要专门的研究。

\std 是啊，现在我想起来了，在\chap{ch-02} 里讨论如何选取摩擦力方向时，我们谈到过这个问题。

\tchr 现在我想更详细地讨论这个问题。我坚信，解考虑摩擦力的习题所遇到的困难，无论对学生还是对某些为物理教材编题的作者来说，显然都估计不足。

\hlred{来看\figr{fig-10} 所示的例子。已知斜面倾角 $\alpha$、物体的重力 $P$、力 $F$ 和摩擦系数（coefficient of friction）$k$。为简单起见，我们假定 $k_{0}=k$，其中 $k_{0}$ 是确定最大可能的静摩擦力的系数。要求确定物体做何种运动，并求出加速度。}
\pnote{原书第 3 章指出：确定最大静摩擦力的系数 $k_{0}$ 略大于滑动摩擦系数 $k$。这与常识相符——使物体起动比维持它滑动要略难一些。本章为分类讨论简便取 $k_{0}=k$；若计及 $k_{0}>k$，下文各个“临界”两侧的静止区间会相应加宽。}

假定物体沿斜面向上滑动。我们可以按\figr{fig-27}~\drkgry{(b)} 所示分解各力，并利用\chap{ch-06} 中得到的加速度结果，于是
\begin{equation}%
a = \frac{g}{P} \left[F \cos \alpha- P \sin \alpha -(P \cos \alpha + F \sin \alpha) \, k \right]
\label{eq-24}
\end{equation}
\pnote{式~\eqref{eq-24} 中的力 $F$ 沿水平方向（\figr{fig-10}；\figr{fig-27} 亦为此情形）：它沿斜面的分量是 $F\cos\alpha$，垂直斜面的分量是 $F\sin\alpha$，故支承力 $N=P\cos\alpha+F\sin\alpha$。代入 $F_{f\!\!r}=kN$，再沿斜面方向用牛顿第二定律，即得该式。}
由式~\eqref{eq-24} 可知，要使物体沿斜面向上滑动，必须满足下列条件：
\begin{equation*}%
F \cos \alpha - P \sin \alpha -\left( P \cos \alpha + F \sin \alpha \right) k \geqslant 0
\end{equation*}
这个条件可以写成
\begin{equation*}%
F \geqslant P \left ( \frac{k \cos \alpha + \sin \alpha}{\cos \alpha -k \sin \alpha} \right)
\end{equation*}
或
\begin{equation}%
F \geqslant P \left ( \frac{k + \tan \alpha}{ 1 -k \tan \alpha} \right)
\label{eq-25}
\end{equation}

我们还要假定斜面的倾角不太大，从而 $(1-k \tan \alpha) > 0$，即
\begin{equation}%
\tan \alpha < \frac{1}{k}
\label{eq-26}
\end{equation}

接下来假定物体沿斜面向下滑动。我们仍按\figr{fig-27}~\drkgry{(b)} 分解各力，但把摩擦力反向。结果得到物体加速度的如下表达式
\begin{equation}%
a= \frac{g}{P} \left[P \sin \alpha -F \cos \alpha - \left( P \cos \alpha + F \sin \alpha \right) \, k \right]
\label{eq-27}
\end{equation}
由式~\eqref{eq-27} 可知，要使物体沿斜面向下滑动，必须满足条件：
\begin{equation*}%
P \sin \alpha -F \cos \alpha - \left( P \cos \alpha + F \sin \alpha \right) k \geqslant 0
\end{equation*}
把这个条件写成
\begin{equation*}%
F \leqslant P \left ( \frac{ \sin \alpha - k \cos \alpha}{\cos \alpha + k \sin \alpha} \right)
\end{equation*}
或
\begin{equation}%
F \leqslant P \left ( \frac{ \tan \alpha -k }{ 1 + k \tan \alpha} \right)
\label{eq-28}
\end{equation}
在这种情形，我们要假定斜面的倾角不太小，从而 $(\tan \alpha- k)>0$，即
\begin{equation}%
\tan \alpha > k
\label{eq-29}
\end{equation}

把条件~\eqref{eq-25}、\eqref{eq-26}、\eqref{eq-28} 与~\eqref{eq-29} 结合起来，可以得出以下结论：
\begin{enumerate}[label=\textbf{(\arabic*)},leftmargin=1cm]
\item 设斜面满足条件
\begin{equation*}%
k < \tan \alpha < \frac{1}{k}
\end{equation*}
那么：
\begin{enumerate}[label=\textbf{(\alph*)},leftmargin=1cm]
\item 若 $F>P \left(\dfrac{k+\tan \alpha}{ 1-k \tan \alpha}\right)$，物体向上滑动，加速度可由式~\eqref{eq-24} 确定；
\item 若 $F=P \left(\dfrac{k+\tan \alpha}{ 1-k \tan \alpha}\right)$，物体以匀速向上滑动，或者静止；
\item 若 $F<P \left(\dfrac{\tan \alpha - k}{ 1+k \tan \alpha}\right)$，物体向下滑动，加速度可由式~\eqref{eq-27} 确定；
\item 若 $F=P \left(\dfrac{\tan \alpha-k}{1+k \tan \alpha}\right)$，物体以匀速向下滑动，或者静止；
\item 若 $P\left( \dfrac{\tan \alpha-k}{1+k \tan \alpha}\right) < F < P \left( \dfrac{k+\tan \alpha}{1-k \tan \alpha} \right)$，物体静止。
\end{enumerate}
注意，当力 $F$ 从 $P \left(\dfrac{\tan \alpha-k}{1 +k \tan \alpha}\right)$ 增大到 $P \left( \dfrac{k+\tan \alpha} {1-k \tan \alpha} \right)$ 时，静摩擦力先从 $k (P \cos \alpha +F \sin \alpha)$ 逐渐减小到零；随后方向反转，再逐渐增大到 $k (P \cos \alpha +F \sin \alpha)$。在此期间物体始终保持静止。
\item 再设斜面满足条件
\begin{equation*}%
0 < \tan \alpha \leqslant k
\end{equation*}
那么：
\begin{enumerate}[label=\textbf{(\alph*)},leftmargin=1cm]
\item 若 $F>P \left(\dfrac{k+\tan \alpha}{ 1-k \tan \alpha}\right)$，物体向上滑动，加速度可由式~\eqref{eq-24} 确定；
\item 若 $F=P \left(\dfrac{k+\tan \alpha}{ 1-k \tan \alpha}\right)$，物体以匀速向上滑动，或者静止；
\item 若 $F<P \left(\dfrac{k+\tan \alpha}{ 1-k \tan \alpha}\right)$，物体静止；物体不可能沿斜面向下滑动（即使力 $F$ 变为零也不可能）。
\end{enumerate}
\item 最后，设斜面满足条件
\begin{equation*}%
\tan \alpha \geqslant \frac{1}{k}
\end{equation*}
那么：
\begin{enumerate}[label=\textbf{(\alph*)},leftmargin=1cm]
\item 若 $F< P \left(\dfrac{\tan \alpha - k}{ 1 + k \tan \alpha}\right)$，物体向下滑动，加速度可由式~\eqref{eq-27} 确定；
\item 若 $F = P \left(\dfrac{\tan \alpha - k}{ 1 + k \tan \alpha}\right)$，物体以匀速向下滑动，或者静止；
\item 若 $F > P \left(\dfrac{\tan \alpha - k}{ 1 + k \tan \alpha}\right)$，物体静止；物体不可能沿斜面向上滑动。乍看起来这似乎不可理解——力 $F$ 不是可以无限增大吗！然而斜面的倾角实在太大了：随着力 $F$ 的增大，物体压紧斜面的压力会以更快的速率增大。
\end{enumerate}
\end{enumerate}
\pnote{结论 2(c) 与 3(c) 的现象，工程上称为自锁（self-locking）。$\tan\alpha\leqslant k$ 即斜面的自锁条件：倾角不大时，物体不受外力也不会下滑；3(c) 则表明，倾角足够大时再大的水平推力也不能使物体上滑。“自锁”一词不在中学课程之内，属物理竞赛与机械常识。}

\std 学校里从来没有给我们讲过这类东西。

\tchr 这正是我想让你注意这个问题的原因。当然，在入学考试中，你们要对付的显然会是简单得多的情形：要么没有摩擦，要么虽有摩擦但运动的性质事先已知（例如已知物体是在运动还是静止）。不过，即使你用不着游过深水处，知道深水处在哪儿也终究是好事。
\pnote{本章这种“先判定运动状态、再列方程”的完整分类讨论超出中学课程要求：人教版必修第一册第三章的摩擦力习题，运动状态大都是给定的（静止、匀速或加速度方向已知）。运动状态未知的摩擦问题在高考中偶见于较难的题目，在物理竞赛中则是常规内容。}

\std 要是假定 $k=0$，会怎么样呢？

\tchr 没有摩擦时，一切立刻变得简单多了。对斜面的任意倾角，结果将是：
\begin{enumerate}[label=$\blacktriangleright$,leftmargin=1cm]
\item 当 $F > P \tan \alpha$ 时，物体向上滑动，加速度为
\begin{equation}%
a = \frac{g}{P} \left( F \cos \alpha -P \sin \alpha \right)
\label{eq-30}
\end{equation}
\item 当 $F = P \tan \alpha$ 时，物体匀速滑动（向上或向下），或者静止；
\item 当 $F < P \tan \alpha$ 时，物体向下滑动，加速度为
\begin{equation}%
a = \frac{g}{P} \left( P \sin \alpha - F \cos \alpha \right)
\label{eq-31}
\end{equation}
\end{enumerate}
注意，式~\eqref{eq-30} 与式~\eqref{eq-31} 的结果只差一个正负号。因此，解题时\hlblue{你尽可以任设一个运动方向，求出 $a$，再注意加速度的符号：若 $a>0$，物体沿你所设的方向运动；若 $a<0$，物体将沿相反的方向运动}（加速度的大小将等于 $|a|$）。
\pnote{这就是中学教辅所说的“假设法”：任设一个方向，按该方向列方程，用所得加速度（或待求力）的符号检验假设。注意它的成立前提——两种相反假设下的方程只差一个正负号。无摩擦时这一前提成立（式~\eqref{eq-30} 与式~\eqref{eq-31}）；有摩擦时不成立：摩擦在两种假设下都阻碍运动，两个方程可能同时给出负的加速度，对应本章开头所分析的“物体静止”，此时符号法失效。}

再考虑一个问题。\hlred{两个物体 $P_{1}$ 和 $P_{2}$ 用跨过滑轮的绳相连。物体 $P_{1}$ 位于倾角为 $\alpha$、摩擦系数为 $k$ 的斜面上；物体 $P_{2}$ 悬挂在绳上\drkgry{（\figr{fig-32}）}。求系统的加速度。}

\begin{figure}[!ht]
\centering
\includegraphics[width=0.4\linewidth]{figs/ch-07/fig-032a.pdf}
\caption[经绳相连的两物体，斜面带摩擦。]{两个物体经跨过滑轮的绳相连，其中一个位于带摩擦的斜面上。问题是求系统的加速度。\label{fig-32}}
\end{figure}

设系统自左向右运动。把系统的运动作为整体来考虑，可以写出加速度的如下方程：
\begin{equation}%
a = g \left( \frac{P_{2} - P_{1} \sin \alpha - P_{1} k \cos \alpha}{P_{1}+P_{2}}\right)
\label{eq-32}
\end{equation}
\pnote{原书以重力的大小 $P_{1}$、$P_{2}$ 标记两物体，故式中出现 $g/(P_{1}+P_{2})$——它就是 $1/(m_{1}+m_{2})$；分子 $P_{2}-P_{1}\sin\alpha-P_{1}k\cos\alpha$ 即系统的合外力（$P_{2}$ 的驱动减去 $P_{1}$ 的下滑分量与摩擦力）。}
再设系统自右向左运动，得到
\begin{equation}%
a = g \left( \frac{P_{1} \sin \alpha - P_{2} - P_{1} k \cos \alpha}{P_{1}+P_{2}}\right)
\label{eq-33}
\end{equation}
现在对给定的 $\alpha$ 和 $k$ 值进行考察，改变比值 $p=P_{2}/P_{1}$。由式~\eqref{eq-32} 可知，要自左向右运动，应当满足条件
\begin{equation*}%
p \leqslant \frac{1}{\sin \alpha + k \cos \alpha}
\end{equation*}
式~\eqref{eq-33} 则表明，要自右向左运动，必要条件是
\begin{equation*}%
p \geqslant \frac{1}{\sin \alpha - k \cos \alpha}
\end{equation*}
这里还需要一个附加条件：倾角不应太小，即 $\tan \alpha >k$。若 $\tan \alpha \leqslant k$，则无论比值 $p$ 多大，系统都不会自右向左运动。

若 $\tan \alpha >k$，则当下列不等式成立时，物体静止：
\begin{equation*}%
 \frac{1}{\sin \alpha + k \cos \alpha} < p < \frac{1}{\sin \alpha - k \cos \alpha}
\end{equation*}
反之，若 $\tan \alpha \leqslant k$，则当
\begin{equation*}%
p > \frac{1}{\sin \alpha + k \cos \alpha}
\end{equation*}
时物体静止。
\pnote{这段条件与正文自己给出的定义 $p=P_{2}/P_{1}$ 不一致。按该定义，由式~\eqref{eq-32}，自左向右运动要求 $p\geqslant\sin\alpha+k\cos\alpha$；由式~\eqref{eq-33}，自右向左要求 $p\leqslant\sin\alpha-k\cos\alpha$（且 $\tan\alpha>k$）；静止区间为 $\sin\alpha-k\cos\alpha<p<\sin\alpha+k\cos\alpha$；若 $\tan\alpha\leqslant k$，则 $p<\sin\alpha+k\cos\alpha$ 即静止。原书印出的四个含 $1/(\sin\alpha\pm k\cos\alpha)$ 的不等式，恰是把 $p$ 换读作 $P_{1}/P_{2}$ 时的正确条件；而正文定义及习题 14 的说明（$p=1/2$ 表示斜面上的物体是悬挂物体的两倍重）都明确 $p=P_{2}/P_{1}$，二者不能同时成立。经核对 1973 年 Mir 英文原版，各处即如此排印，当为原书（或其英译）自身的疏漏而非录入之误；译文照原书排印，谨此说明。另：底本本章另有若干拼写级排印小误（如 ``fall owing''、``down ward''、``in, solving''），径改不注。}

\std 那要是改变角 $\alpha$ 或系数 $k$，会怎么样呢？

\tchr 从这个角度进行的考察，我留给你作为家庭作业（见习题 13 和 14）。

\end{dialogue}
\section*{习题}
\label{prob-07}
\addcontentsline{toc}{section}{习题}

\begin{enumerate}[resume*=problems]
\item 对\figr{fig-32} 所示的问题进行考察：设斜面倾角 $\alpha$ 与比值 $p=P_{2}/P_{1}$ 给定，给摩擦系数 $k$ 取各种不同的值。
\item 对\figr{fig-32} 所示的问题进行考察：设摩擦系数 $k$ 与比值 $p=P_{2}/P_{1}$ 给定，给斜面倾角 $\alpha$ 取各种不同的值。为简单起见，比值只取两个值：$p= 1$（两物体等重）和 $p= 1/2$（斜面上的物体是悬挂在绳上的物体的两倍重）。
\end{enumerate}
\cleardoublepage

\thispagestyle{empty}
\vspace*{-2cm}
\begin{figure*}[!ht]
\centering
\includegraphics[width=0.68\textwidth]{figs/sec/tarasov-qasp-sec-04.pdf}
\end{figure*}
\begin{fullwidth}
{\textsf{\Large \hlred{圆周运动是曲线运动最简单的形式，正因如此，领会这种运动的特点也就更加重要。你可以看到，整个宇宙无不由曲线运动构成。让我们来讨论质点沿圆周的匀速与非匀速运动，以及卫星的轨道运动——这将把我们引向对物体失重现象物理成因的讨论。}}}
\end{fullwidth}
% NEWTERMS: force of static friction | 静摩擦力 | 已由第 3 章首现括注（static friction），本章沿用不重复括注；k0 见第 3 章
% NEWTERMS: coefficient of friction | 摩擦系数 | 泛称；滑动摩擦系数已在第 2 章定名
% NEWTERMS: uniform velocity / uniform motion | 匀速（运动） | 第 8 章起用于“匀速圆周运动”
% NEWTERMS: curvilinear motion | 曲线运动 | 章尾格言首现
% NEWTERMS: weightlessness | 失重 | 章尾格言首现，第 8–9 章主题
% NEWTERMS: orbiting satellite | （沿轨道运行的）卫星 | 章尾格言
% NEWTERMS: self-locking | 自锁 | 仅评注使用；竞赛与机械常识概念
% NEWTERMS: home assignment | 家庭作业 | 对应章末习题（全书连续编号，本章为 13、14）
